\section{Keeping the exact marginal probabilities}
\label{p13:exact:section}

The optimal remainder is smaller than \(N^{-1}\), the cost of deleting
all irregular windows. We therefore keep every event with nonzero
probability from the outset. Its effective square-class space records
all local relations once, and its actual probability is the correct
centering parameter.

Let \(\mathcal I_N=\{x\in I_N:p_x>0\}\), where \(p_x=\E J_x\).
An event outside \(\mathcal I_N\) is identically zero and can be
removed without changing \(Z_N\). For every \(x\in I_N\), let
\(S_x\) be the image of the even selections of
\(V_x=\{x-1,\ldots,x+L-1\}\) in
\(\mathbb Q^\times/\mathbb Q^{\times2}\), regarded as a vector
space over \(\F_2\). Write \(d_x=\dim S_x\).
The even-selection space has dimension $L$. Its kernel consists of
local even square relations. The following bound makes the possible
effective ranks explicit.

We apply the binary Griesmer bound
\cite[Section~2, equation~(2.10)]{griesmer1960} to the code of
local even square relations. The arithmetic input is its minimum
nonzero weight; the residual-code proof below keeps the dependence
on the parameters explicit.

\begin{lemma}[Local rigidity]
\label{p1:limits:local-rigidity}
Let $V=\{U,\ldots,U+B-1\}$, where $U\asymp N$ and $B=O(\log N)$.
Put
\[
 d_{N,B}=\left\lceil\frac{\log U}{\log(64B)}\right\rceil.
\]
If the relative square-relation code has dimension $r\ge t$, then
\begin{equation}
 \pi(B)+1\ge
 \sum_{i=0}^{t-1}\left\lceil d_{N,B}/2^i\right\rceil-t.
 \label{p1:limits:griesmer}
\end{equation}
Consequently $r\le1$ when $B=\log N/\log2+o(\log N)$.
More generally, if $B\sim c\log N$ and
$c<2-2^{1-t}$, then $r<t$ for all sufficiently large $N$.
For every fixed $\varepsilon>0$, the relative signs are exactly
independent in all windows with $B\le(1-\varepsilon)\log N$.
\end{lemma}
\begin{proof}
Every square relation is supported on the $D$ integers whose odd
prime valuations are all at primes at most $B$: a larger prime is
private in the window.  The even square-relation code on these
coordinates has redundancy at most $\pi(B)+1$.  Its minimum nonzero
weight is at least $d_{N,B}$ by Lemma~\ref{p1:arith:runge}.

For completeness, a binary linear code of length $D$, dimension $r$
and minimum distance $d$ has redundancy at least
$\sum_{i<t}\lceil d/2^i\rceil-t$ whenever $r\ge t$.
Choose a word of minimum weight $d_0\ge d$ and project outside its
support.  The kernel of this projection is its one-dimensional span:
another nonzero word supported there would have smaller weight.
The projected code has dimension $r-1$ and minimum distance at least
$\lceil d_0/2\rceil$, since one of two complementary lifts uses at
most $\lfloor d_0/2\rfloor$ coordinates on that support.  The
redundancy decreases by $d_0-1$.  Repeating $t$ times proves the
claim, and hence \eqref{p1:limits:griesmer}.

The prime number theorem gives
$\pi(B)\sim c\log N/\log\log N$, whereas
$d_{N,B}\sim\log N/\log\log N$ when $B\sim c\log N$.
Thus \eqref{p1:limits:griesmer} excludes dimension $t$ under the
stated strict inequality.  At critical length, $1/\log2<3/2$,
so $t=2$ applies.  The assertion for every smaller subcritical
window follows by including it in a window of length
$\lfloor(1-\varepsilon/2)\log N\rfloor$ and projecting.
\end{proof}

For an admissible window the affine prescription is consistent. Its
probability is $2^{-d_x}$, as the character expansion in
Section~\ref{p1:alg:section} also shows. The lemma therefore gives
\begin{equation}
 d_x\in\{L-1,L\},\qquad p_x=2^{-d_x}\in\{p,2p\},
 \qquad p_*:=\max_{x\in\mathcal I_N}p_x\ll_C N^{-1}.
 \label{p13:exact:marginals}
\end{equation}
The assertion about \(d_x\) holds for all sufficiently large \(N\)
in a fixed critical band. A local square relation may instead
contradict the run prescription; that is precisely the case
already removed by \(p_x=0\).

For a finite \(I\subseteq\mathcal I_N\), put
\begin{align}
 \rho_{\mathrm{rel}}(I)
  &=\sum_{x\in I}d_x-\dim\sum_{x\in I}S_x,
       & d_{xy}&=\dim(S_x\cap S_y),\label{p13:exact:rank}\\
 q(I)&=\E\prod_{x\in I}J_x,
       & w(I)&=2^{\rho_{\mathrm{rel}}(I)}
                         \prod_{x\in I}p_x,\label{p13:exact:weights}\\
 c(I)&=\E\prod_{x\in I}(J_x-p_x).
 \label{p13:exact:center}
\end{align}
The empty products are one. We have \(q(I)=w(I)\) when the
affine prescriptions are jointly compatible, and \(q(I)=0\)
otherwise. On the regular set \(G_N\), these definitions agree
with the homogeneous-marginal notation used in the companion.
Outside \(G_N\), \(\rho_{\mathrm{rel}}\) counts dependencies
between effective spaces; it does not count local relations a
second time.

The notions of directness, connectedness, and primitive family
refer to the direct-sum decomposition of the effective spaces
proved in Section~\ref{p1:alg:section}. A family is direct when its
effective spaces are in direct sum, and connected when it has just one
canonical direct component. In particular, a primitive
family has a nonzero relation and every proper subfamily is
direct. Every nonzero relation of a primitive family has a
nonzero projection in each of its blocks. A connected family
need not have a relation with this property on its entire
support. This distinction matters in the conditional count below.

For fixed \(r\ge2\), define
\begin{equation}
 A_r=\sum_{\substack{I\subseteq\mathcal I_N\\|I|=r}}|c(I)|,
 \qquad
 K_3^+=\sum_{\substack{|T|=3,\ T\text{ connected}\\q(T)>0}}c(T).
 \label{p13:exact:activities}
\end{equation}
The nonnegativity in the definition of \(K_3^+\) will be proved
in Proposition~\ref{p13:activity:kappa3}.
The exact second and third factorial cumulants satisfy
\begin{align}
 \alpha_N&=2\sum_{|e|=2}c(e)-\sum_xp_x^2,
 \label{p13:exact:alpha}\\
 \kappa_{3,N}
   &=6\sum_{|T|=3}c(T)
       -6\sum_{e=\{x,y\}}(p_x+p_y)c(e)
       +2\sum_xp_x^3.
 \label{p13:exact:kappa3}
\end{align}
These identities follow by expanding the first three factorial
moments. They display the repeated-index terms that an
independent Bernoulli reference will retain exactly.

\paragraph{Fixed sources on effective spaces.}
The fixed-source estimates stated in
Propositions~\ref{p1:fixed:targets} and~\ref{p1:fixed:propagation}
apply on
\(\mathcal I_N\), with the new conditional rank
\[
 \delta_z=\dim\bigl(S_z\cap\textstyle\sum_{x\in I}S_x\bigr)
\]
and weight \(p_z2^{\delta_z}\). To see this, choose in each
even-selection space a linear complement to its local kernel.
The valuation map is injective on that complement and maps it
onto \(S_x\). Every label regularization, prime row, and disjoint
square equation in the proofs is applied to this linear
realization. Restriction to a complement restricts the allowed
binary selections within the same physical coordinates; it does
not enlarge any prime support.
The factor \(p_z\) is at most \(2p\), and the number of sources
and targets is fixed in those propositions. Thus their constants
change by a fixed factor. The same reasoning applies to
homogeneous prescriptions. In particular, for fixed \(s,m\) and
\(\varepsilon>0\),
\begin{equation}
 \sum_{\substack{T\subseteq\mathcal I_N\setminus I\\|T|=m}}
       q(I\cup T)\ll_{s,m,C,\varepsilon}N^\varepsilon q(I)
 \quad (|I|=s),
 \label{p13:exact:propagation}
\end{equation}
and the analogous homogeneous bound holds with \(q\) replaced
by \(w\). If \(q(I)=0\), every summand on the left is zero.
No estimate here is asserted for \(s\) or \(m\) growing with \(N\).
